\section{Fibonacci windows and the all-\texorpdfstring{$q$}{q} moment transfer}
\label{sec:second-collision}

By \Cref{thm:partition-difference}, the finite-window fiber spectrum is already
encoded by the classical Fibonacci partition function.  The collision moments
are therefore naturally expressed as power sums of a Fibonacci-lag discrete
derivative.  The key point of this section is that those power sums are
precisely trapped between two adjacent Fibonacci windows of the classical
partition power sums.  This closes the finite-window moment problem to an
all-\(q\) theorem chain.

\subsection{Difference-power sums and Fibonacci windows}

\begin{definition}[Difference spectrum]
\label{def:difference-spectrum}
For \(0\le n<F_{m+2}\), define
\[
  a_m(n):=d_m^\#(\pi_m(n))
  =
  R^\dagger(n)-R^\dagger(n-F_{m+1}).
\]
For \(q\ge 1\), define
\[
  S_q(m):=\sum_{x\in X_m} d_m(x)^q,
  \qquad
  T_q^\dagger(N):=\sum_{n=0}^{N-1} R^\dagger(n)^q.
\]
\end{definition}

\begin{theorem}[Difference-power sums]
\label{thm:difference-power-sums}
For every integer \(q\ge 1\),
\[
  S_q(m)
  =
  \sum_{n=0}^{F_{m+2}-1} a_m(n)^q
  =
  \sum_{n=0}^{F_{m+2}-1}
  \bigl(R^\dagger(n)-R^\dagger(n-F_{m+1})\bigr)^q.
\]
\end{theorem}

\begin{proof}
Because \(\pi_m\) is a permutation of \(\ZZ/F_{m+2}\ZZ\),
\[
  S_q(m)
  =
  \sum_{r=0}^{F_{m+2}-1} d_m^\#(r)^q
  =
  \sum_{n=0}^{F_{m+2}-1} d_m^\#(\pi_m(n))^q.
\]
Now apply \Cref{thm:partition-difference}.
\end{proof}

\begin{theorem}[Fibonacci-window sandwich]
\label{thm:fibonacci-window-sandwich}
For every pair of integers \(q\ge 1\) and \(m\ge 1\),
\[
  T_q^\dagger(F_{m+1})\le S_q(m)\le T_q^\dagger(F_{m+2}).
\]
\end{theorem}

\begin{proof}
By \Cref{thm:difference-power-sums}, for \(0\le n<F_{m+1}\) we have
\(R^\dagger(n-F_{m+1})=0\), hence \(a_m(n)=R^\dagger(n)\).  Therefore
\[
  \sum_{n=0}^{F_{m+1}-1} a_m(n)^q
  =
  \sum_{n=0}^{F_{m+1}-1} R^\dagger(n)^q
  =
  T_q^\dagger(F_{m+1}),
\]
which gives the lower bound.

For the upper bound, \(a_m(n)\) is a fiber multiplicity and hence nonnegative.
By definition,
\[
  a_m(n)=R^\dagger(n)-R^\dagger(n-F_{m+1})\le R^\dagger(n)
  \qquad (0\le n<F_{m+2}).
\]
Raise both sides to the \(q\)-th power and sum over \(n\).
\end{proof}

\begin{theorem}[Transfer of partition constants]
\label{thm:all-q-transfer}
For every integer \(q\ge 1\),
\[
  T_q^\dagger(N)\asymp_q N^{(\log \lambda_q)/\log\varphi},
\]
where \(\lambda_q\) is Sanna's partition constant \cite[Thm.~1.1]{Sanna2025}.  Hence
\[
  S_q(m)\asymp_q \lambda_q^m.
\]
If \(q\ge 2\), then the Perron root \(r_q\) of the finite-state collision
kernel of \Cref{thm:collision-kernel} satisfies
\[
  r_q=\lambda_q.
\]
\end{theorem}

\begin{proof}
Pointwise,
\[
  R(n)\le R^\dagger(n)=R(n)+R(n-1)\le 2\max\{R(n),R(n-1)\},
\]
so
\[
  R(n)^q\le R^\dagger(n)^q
  \le 2^{q-1}\bigl(R(n)^q+R(n-1)^q\bigr).
\]
Summing over \(0\le n<N\) yields
\[
  \sum_{n=0}^{N-1} R(n)^q
  \le
  T_q^\dagger(N)
  \le
  2^q\sum_{n=0}^{N-1} R(n)^q.
\]
Sanna's theorem \cite[Thm.~1.1]{Sanna2025} therefore gives
\[
  T_q^\dagger(N)\asymp_q N^{(\log \lambda_q)/\log\varphi}.
\]
Taking \(N=F_{m+1}\) and \(N=F_{m+2}\), using \(F_m\asymp\varphi^m\), and
applying \Cref{thm:fibonacci-window-sandwich}, we obtain
\[
  S_q(m)\asymp_q \lambda_q^m.
\]
If \(q\ge 2\), let \(A_q,u_q,v_q\) be the trimmed realization from
\Cref{thm:collision-kernel}, and write \(r_q:=\rho(A_q)\).  For every
\(\varepsilon>0\), Gelfand's formula yields
\(\|A_q^m\|\ll_{\varepsilon,q}(r_q+\varepsilon)^m\), and therefore
\[
  S_q(m)=u_q^{\mathsf T}A_q^m v_q\ll_{\varepsilon,q}(r_q+\varepsilon)^m.
\]
Since \(S_q(m)\asymp_q \lambda_q^m\), this implies
\(\lambda_q\le r_q+\varepsilon\) for every \(\varepsilon>0\), hence
\(\lambda_q\le r_q\).

For the reverse inequality, choose a strongly connected component \(B\) of
\(A_q\) with \(\rho(B)=r_q\).  Because the automaton is trimmed, there exist
states \(s,t\) of \(B\) and integers \(a,b\ge 0\) such that the initial state
reaches \(s\) in \(a\) steps and \(t\) reaches an accepting state in \(b\)
steps.  Hence
\[
  S_q(m+a+b)\ge (B^m)_{st}
  \qquad (m\ge 0).
\]
Since \(B\) is irreducible, Perron--Frobenius theory gives
\[
  \limsup_{m\to\infty}(B^m)_{st}^{1/m}=r_q
\]
\cite[Ch.~1]{Seneta2006}.  On the other hand,
\(S_q(m+a+b)\asymp_q \lambda_q^{m+a+b}\), so
\[
  \lambda_q
  =
  \lim_{m\to\infty}S_q(m+a+b)^{1/m}
  \ge
  \limsup_{m\to\infty}(B^m)_{st}^{1/m}
  =
  r_q.
\]
Thus \(r_q=\lambda_q\).
\end{proof}

\begin{remark}[The case \texorpdfstring{$q=1$}{q=1}]
\label{rem:lambda-one}
For \(q=1\), \(S_1(m)=2^m\) gives \(\lambda_1=2\) directly.
\end{remark}

\begin{remark}[Relation to classical moment problems]
\label{rem:sanna-relation}
Theorem~\ref{thm:difference-power-sums} places the finite-window collision
moments beside the ordinary power sums of the Fibonacci partition function
studied by Sanna \cite{Sanna2025}.  Theorem~\ref{thm:fibonacci-window-sandwich}
sharpens that relation: every finite-window collision moment is exactly trapped
between two adjacent Fibonacci windows of the indexed partition power sums.
Theorem~\ref{thm:all-q-transfer} then shows that the fold moments share the
same exponential constants \(\lambda_q\).
\end{remark}

\begin{corollary}[Kernel formula]
\label{cor:multivariate-moment-kernel}
Fix \(q\ge 2\), write \(\mathbf z=(z_1,\dots,z_{q-1})\), and for
\(A\subseteq [q-1]\) write
\[
  z_A:=\prod_{i\in A}z_i,
  \qquad
  z_\varnothing:=1.
\]
Then
\[
  S_q(m)
  =
  [\mathbf z^0]
  \prod_{j=1}^m
  \left(
    2+\sum_{\varnothing\neq A\subseteq [q-1]}
    \bigl(z_A^{F_j}+z_A^{-F_j}\bigr)
  \right).
\]
\end{corollary}

\begin{proof}
By \Cref{thm:affine-transfer},
\[
  S_q(m)
  =
  \sum_{n=0}^{F_{m+2}-1}\widetilde d_m(n)^q.
\]
Let
\[
  \mathcal P_m(z):=\prod_{j=1}^m(1+z^{F_j})
  =
  \sum_{n=0}^{F_{m+2}-1}\widetilde d_m(n)z^n.
\]
Diagonal extraction gives
\[
  S_q(m)
  =
  [\mathbf z^0]\,
  \mathcal P_m(z_1)\cdots \mathcal P_m(z_{q-1})
  \mathcal P_m\bigl((z_1\cdots z_{q-1})^{-1}\bigr).
\]
Using the product form of \(\mathcal P_m\), we obtain
\[
  S_q(m)
  =
  [\mathbf z^0]
  \prod_{j=1}^m
  \left(
    \prod_{i=1}^{q-1}(1+z_i^{F_j})
  \right)
  \left(
    1+(z_1\cdots z_{q-1})^{-F_j}
  \right).
\]
Finally,
\begin{align*}
  \prod_{i=1}^{q-1}(1+z_i)\left(1+(z_1\cdots z_{q-1})^{-1}\right)
  &=
  \sum_{A\subseteq [q-1]}z_A
  +
  \sum_{A\subseteq [q-1]}z_{A^c}^{-1} \\
  &=
  2+\sum_{\varnothing\neq A\subseteq [q-1]}(z_A+z_A^{-1}),
\end{align*}
and replacing \(z_i\) by \(z_i^{F_j}\) yields the claim.
\end{proof}

\subsection{The quadratic case}

\begin{definition}[Quadratic difference energy]
\label{def:T2}
Set
\[
  T_2(m):=S_2(m)=\sum_{n=0}^{F_{m+2}-1} a_m(n)^2.
\]
\end{definition}

\begin{theorem}[Cubic recurrence]
\label{thm:cubic-recurrence}
The quadratic difference energy satisfies
\[
  T_2(1)=2,\qquad T_2(2)=6,\qquad T_2(3)=14,
\]
and for every \(m\ge 4\),
\[
  T_2(m)=2T_2(m-1)+2T_2(m-2)-2T_2(m-3).
\]
\end{theorem}

\begin{proof}
By \Cref{thm:partition-difference},
\[
  T_2(m)
  =
  \sum_{n=0}^{F_{m+2}-1}d_m^\#(\pi_m(n))^2
  =
  \sum_{n=0}^{F_{m+2}-1}\widetilde d_m(n)^2.
\]
Let
\[
  \mathcal P_m(z):=\prod_{j=1}^m(1+z^{F_j})
  =
  \sum_{n=0}^{F_{m+2}-1}\widetilde d_m(n)z^n.
\]
Then
\[
  T_2(m)
  =
  [z^0]\bigl(\mathcal P_m(z)\mathcal P_m(z^{-1})\bigr).
\]
Define
\[
  B_m(t):=[z^t]\prod_{j=1}^m\bigl(2+z^{F_j}+z^{-F_j}\bigr),
  \qquad t\in\ZZ.
\]
Then \(B_m(t)=B_m(-t)\) and \(T_2(m)=B_m(0)\).

The defining product yields the coefficient recursion
\begin{equation}\label{eq:Bm-rec}
  B_m(t)=2B_{m-1}(t)+B_{m-1}(t-F_m)+B_{m-1}(t+F_m).
\end{equation}
Taking \(t=0\) and using symmetry gives
\begin{equation}\label{eq:T-U}
  T_2(m)=2T_2(m-1)+2B_{m-1}(F_m).
\end{equation}
Introduce the auxiliary quantity
\[
  U_m:=B_m(F_{m+1}).
\]
Then \eqref{eq:T-U} becomes
\begin{equation}\label{eq:T-U2}
  T_2(m)=2T_2(m-1)+2U_{m-1}.
\end{equation}

Apply \eqref{eq:Bm-rec} with \(t=F_{m+1}\):
\[
  U_m
  =
  2B_{m-1}(F_{m+1})
  +B_{m-1}(F_{m+1}-F_m)
  +B_{m-1}(F_{m+1}+F_m).
\]
The support of \(\prod_{j=1}^{m-1}(2+z^{F_j}+z^{-F_j})\) lies in
\[
  |t|\le \sum_{j=1}^{m-1}F_j=F_{m+1}-1,
\]
so
\[
  B_{m-1}(F_{m+1})=0
  \qquad\text{and}\qquad
  B_{m-1}(F_{m+1}+F_m)=0.
\]
Since \(F_{m+1}-F_m=F_{m-1}\), we obtain
\begin{equation}\label{eq:U-shift}
  U_m=B_{m-1}(F_{m-1}).
\end{equation}

Now apply \eqref{eq:Bm-rec} with index \(m-1\) and \(t=F_{m-1}\):
\[
  B_{m-1}(F_{m-1})
  =
  2B_{m-2}(F_{m-1})
  +B_{m-2}(0)
  +B_{m-2}(2F_{m-1}).
\]
Again the support bound
\[
  \sum_{j=1}^{m-2}F_j=F_m-1<2F_{m-1}
\]
forces \(B_{m-2}(2F_{m-1})=0\).  Therefore
\begin{equation}\label{eq:U-rec}
  U_m=2U_{m-2}+T_2(m-2).
\end{equation}
Replacing \(m\) by \(m-1\) in \eqref{eq:U-rec} gives
\[
  U_{m-1}=2U_{m-3}+T_2(m-3).
\]
From \eqref{eq:T-U2} with index \(m-2\),
\[
  T_2(m-2)=2T_2(m-3)+2U_{m-3},
\]
so
\begin{equation}\label{eq:U-elim}
  2U_{m-1}=2T_2(m-2)-2T_2(m-3).
\end{equation}
Substituting \eqref{eq:U-elim} into \eqref{eq:T-U2} yields
\[
  T_2(m)=2T_2(m-1)+2T_2(m-2)-2T_2(m-3).
\]
The initial values are obtained directly from the cases \(m=1,2,3\).
\end{proof}

\begin{corollary}[Generating function]
\label{cor:sigma-generating}
The sequence \((T_2(m))_{m\ge 1}\) has generating function
\[
  \sum_{m\ge 1}T_2(m)u^m
  =
  \frac{2u+2u^2-2u^3}{1-2u-2u^2+2u^3}.
\]
\end{corollary}

\begin{proof}
Apply the standard generating-function calculation to the recurrence of
\Cref{thm:cubic-recurrence}.
\end{proof}

\begin{theorem}[Alternating second-order term]
\label{thm:q2-alternating-second-order}
Let
\[
  p(\lambda):=\lambda^3-2\lambda^2-2\lambda+2.
\]
Then \(p\) has exactly three distinct real roots
\[
  \lambda_-\in(-2,-1),
  \qquad
  \lambda_+\in(0,1),
  \qquad
  \rho\in(2,3).
\]
Consequently there exist real constants \(c_\rho,c_-,c_+\) such that
\[
  T_2(m)=c_\rho\rho^m+c_-\lambda_-^m+c_+\lambda_+^m
  \qquad (m\ge 1),
\]
with \(c_\rho>0\) and \(c_-\neq 0\).  In particular,
\[
  T_2(m)-c_\rho\rho^m
  =
  c_-\lambda_-^m+O(\lambda_+^m),
\]
so the second-order term is eventually alternating in sign.
\end{theorem}

\begin{proof}
The sign checks
\[
  p(-2)=-2,\qquad p(-1)=1,\qquad p(0)=2,
\]
\[
  p(1)=-1,\qquad p(2)=-2,\qquad p(3)=5
\]
show that \(p\) has roots in the three disjoint intervals
\((-2,-1)\), \((0,1)\), and \((2,3)\).  Since \(p\) is cubic, these are all
its roots, and they are distinct.

By \Cref{thm:cubic-recurrence}, \((T_2(m))\) is a linear recurrence with
characteristic polynomial \(p\).  Since the roots are simple, there exist
constants \(c_\rho,c_-,c_+\in\RR\) such that
\[
  T_2(m)=c_\rho\rho^m+c_-\lambda_-^m+c_+\lambda_+^m.
\]
Because \(T_2(m)>0\) and \(\rho\) is the unique dominant root, one has
\(c_\rho>0\).

It remains to show \(c_-\neq 0\).  Suppose instead that \(c_-=0\).  Then
\((T_2(m))\) would satisfy the second-order recurrence with roots \(\rho\) and
\(\lambda_+\):
\[
  T_2(m+2)=(\rho+\lambda_+)T_2(m+1)-\rho\lambda_+T_2(m).
\]
Applying this at \(m=1\) and using \((T_2(1),T_2(2),T_2(3))=(2,6,14)\) yields
\[
  14=6(\rho+\lambda_+)-2\rho\lambda_+.
\]
By Vieta's formulas for \(p\),
\[
  \rho+\lambda_+=2-\lambda_-
  \qquad\text{and}\qquad
  \rho\lambda_+=\frac{-2}{\lambda_-}.
\]
Substituting these identities gives
\[
  14=12-6\lambda_-+\frac{4}{\lambda_-},
\]
so \(\lambda_-\) would satisfy
\[
  3\lambda_-^2+\lambda_--2=0.
\]
Dividing \(p(\lambda)\) by \(3\lambda^2+\lambda-2\) gives remainder
\((4-5\lambda)/9\), so any common root would satisfy \(\lambda=4/5\), which is
impossible because \(\lambda_-\in(-2,-1)\).  This contradiction proves
\(c_-\neq 0\).  Since \(|\lambda_-|>\lambda_+\) and \(\lambda_-<0\), the
second-order term is eventually alternating.
\end{proof}

\begin{remark}[Chow--Jones]
\label{rem:chow-jones}
Chow and Jones \cite{ChowJones2024} showed that the variance of the Fibonacci
partition function is governed by the same cubic
\(\lambda^3-2\lambda^2-2\lambda+2\).  Theorem~\ref{thm:cubic-recurrence} gives
a finite-window explanation: it controls the quadratic energy of the
Fibonacci-lag discrete derivative \(a_m(n)\).  Since
\Cref{thm:all-q-transfer} yields \(T_2(m)\asymp \lambda_2^m\), its dominant
root \(\rho\) is exactly \(\lambda_2\).
\end{remark}
